\section{Entendendo as camadas de uma imagem}

Uma imagem de container, quando criada, é formada por \textbf{camadas imutáveis}. Cada camada representa um conjunto de modificações no sistema de arquivos.

\subsection{Exemplo de camadas}

Considere uma imagem Python típica:

\begin{enumerate}
    \item A primeira camada adiciona comandos básicos e um gerenciador de pacotes (\texttt{apt}).
    \item A segunda camada adiciona o runtime Python e o \texttt{pip}.
    \item A terceira camada copia o arquivo \texttt{requirements.txt} da aplicação.
    \item A quarta camada instala as dependências listadas.
    \item A quinta camada copia o código-fonte da aplicação.
\end{enumerate}

Essa estrutura em camadas é benéfica porque permite a \textbf{reutilização} entre imagens. Se duas aplicações utilizam Python, ambas podem compartilhar as camadas que instalam o Python e o \texttt{pip}, tornando os builds mais rápidos e reduzindo o armazenamento necessário.

\subsection{Empilhamento de camadas}

\begin{enumerate}
    \item Cada camada, ao ser instalada, é extraída em seu próprio diretório no sistema de arquivos local.
    \item Quando um container é iniciado, um \textbf{sistema de arquivos unificado} é criado empilhando todas as camadas, gerando uma visão unificada.
    \item O container então redireciona seu diretório raiz para essa visão unificada usando \texttt{chroot}.
\end{enumerate}

\subsection{Criando layers manualmente}

Embora na prática se use Dockerfiles, é possível criar layers manualmente com \texttt{docker container commit} para entender o mecanismo interno.

\subsubsection{Criando uma base image}

\begin{verbatim}
# Iniciar container Ubuntu interativo
docker run --name=base-container -ti ubuntu

# Dentro do container: instalar Node.js
apt update && apt install -y nodejs
node -e 'console.log("Hello world!")'

# Em outro terminal: salvar como nova imagem
docker container commit -m "Add node" base-container node-base

# Verificar as layers
docker image history node-base

# Testar a nova imagem
docker run node-base node -e "console.log('Hello again')"

# Remover o container base
docker rm -f base-container
\end{verbatim}

Uma \textbf{base image} é uma fundação para construir outras imagens. Qualquer imagem pode ser usada como base, mas algumas são criadas intencionalmente como blocos de construção.

\subsubsection{Construindo uma app image}

\begin{verbatim}
# Iniciar container a partir da node-base
docker run --name=app-container -ti node-base

# Dentro do container: criar um programa Node
echo 'console.log("Hello from an app")' > app.js
node app.js

# Em outro terminal: salvar como nova imagem com comando padrão
docker container commit -c "CMD node app.js" -m "Add app" \
    app-container sample-app

# Verificar as layers
docker image history sample-app

# Rodar a imagem final
docker run sample-app

# Limpar
docker rm -f app-container
\end{verbatim}

A flag \texttt{-c} define configurações adicionais na imagem --- neste caso, o comando padrão ao iniciar o container será \texttt{node app.js}.

Na prática, a maioria dos builds não usa \texttt{docker container commit}. Em vez disso, utiliza-se um \textbf{Dockerfile} que automatiza todos esses passos.
